\section{Cost Evaluation and Execution Policies}
\label{sec:cost}

\subsection{Metric Vector}

Every plan carries the five-dimensional metric vector of
Section~\ref{sec:math}: latency ($m_1$, ms), peak memory ($m_2$, MB),
synchronization overhead ($m_3$, ms), energy ($m_4$, mJ), and
numerical error ($m_5$, dimensionless $L_2$). The dimensions are
deliberately heterogeneous: they capture qualitatively different
failure budgets (time, space, coordination, power, precision), and the
weight vector---not the metric definitions---decides their relative
importance.

\subsection{Policies as Weight Vectors}

\begin{table}[H]
\centering
\small
\caption{Preset execution policies.}
\label{tab:policies}
\begin{tabular}{@{}lll@{}}
\toprule
Policy & Weight vector $\mathbf{w}$ & Recommended context \\
\midrule
Latency  & $(0.6, 0.1, 0.1, 0.1, 0.1)$ & real-time, interactive services \\
Memory   & $(0.1, 0.6, 0.1, 0.1, 0.1)$ & edge devices, multi-model batching \\
Energy   & $(0.1, 0.1, 0.1, 0.6, 0.1)$ & mobile, battery, green computing \\
Balanced & $(0.2, 0.2, 0.2, 0.2, 0.2)$ & general-purpose, benchmarking \\
Custom   & user-supplied               & domain-specific requirements \\
\bottomrule
\end{tabular}
\end{table}

Policy resolution is a constant-time lookup plus normalization. By
Theorem~\ref{thm:policy-independence}, switching policies can never
change which plans are \emph{valid}---only which valid plan wins.

\subsection{Normalization}

Because the metric dimensions carry different units, implementations
normalize each dimension before applying the weights; the reference
implementations divide by the per-dimension maximum across the valid
candidate set. Normalization choice is behavior-relevant: as the
empirical study shows (Section~\ref{subsec:empirical}, Finding~2),
max-normalization amplifies small absolute differences in a dimension
whose spread is tiny---a $10^{-6}$ relative numerical deviation
normalizes to $1.0$ when it is the largest in its dimension, which
makes selection strongly exactness-preserving. Threshold-relative
normalization ($m_5 / \epsilon$) is a principled alternative and is
discussed in Section~\ref{sec:discussion}.
